\begin{itemize}[label={},leftmargin=1.27cm,itemindent=-1.27cm,itemsep=0.6\baselineskip]
\tightlist
\item
  Adler, G., Djigbenou, M.-L., y Sosa, S. (2016). Global financial
  shocks and foreign asset repatriation: Do local investors play a
  stabilizing role? \emph{Journal of International Money and Finance},
  \emph{60}, 8-28. https://doi.org/10.1016/j.jimonfin.2015.03.007
\item
  Adler, G., y Sosa, S. (2013). \emph{External conditions and debt
  sustainability in Latin America} (IMF Working Paper No.~13/27). Fondo
  Monetario Internacional.
\item
  Ando, S., Mishra, P., Patel, N., Peralta-Alva, A., y Presbitero, A. F.
  (2025). Fiscal consolidation and public debt. \emph{Journal of
  Economic Dynamics and Control}, \emph{170}, artículo 104998.
  https://doi.org/10.1016/j.jedc.2024.104998
\item
  Aruguete, E. (2021). Crisis de balanza de pagos y recuperación de la
  sostenibilidad de la deuda pública en la Argentina (2016-2020).
  \emph{Revista Estado y Políticas Públicas}, (16), 101-123.
\item
  Bai, J., y Perron, P. (2003). Computation and analysis of multiple
  structural change models. \emph{Journal of Applied Econometrics},
  \emph{18}(1), 1-22. https://doi.org/10.1002/jae.659
\item
  Banco Central de la República Argentina {[}BCRA{]}. (s. f.).
  \emph{Balance del Banco Central, serie histórica} {[}Conjunto de
  datos{]}.
  https://www.bcra.gob.ar/Pdfs/PublicacionesEstadisticas/balbcrhis.xls
\item
  Banco Central de la República Argentina {[}BCRA{]}. (2025).
  \emph{Informe monetario mensual, junio de 2025}.
  https://www.bcra.gob.ar/archivos/Pdfs/PublicacionesEstadisticas/informe-monetario-mensual-jun-25.pdf
\item
  Banco Central de la República Argentina {[}BCRA{]}. (2026).
  \emph{Informe monetario mensual, diciembre de 2025}.
  https://www.bcra.gob.ar/archivos/Pdfs/PublicacionesEstadisticas/informe-monetario-mensual-dic-25.pdf
\item
  Blanchard, O. (1990). \emph{Suggestions for a new set of fiscal
  indicators} (OECD Economics Department Working Papers No.~79). OECD
  Publishing.
\item
  Blanchard, O. (2019). Public debt and low interest rates.
  \emph{American Economic Review}, \emph{109}(4), 1197-1229.
  https://doi.org/10.1257/aer.109.4.1197
\item
  Blanchard, O., y Perotti, R. (2002). An empirical characterization of
  the dynamic effects of changes in government spending and taxes on
  output. \emph{Quarterly Journal of Economics}, \emph{117}(4),
  1329-1368. https://doi.org/10.1162/003355302320935043
\item
  Board of Governors of the Federal Reserve System (US). (s. f.).
  \emph{Federal funds effective rate} {[}FEDFUNDS; conjunto de datos{]}.
  Federal Reserve Bank of St.~Louis, FRED.
  https://fred.stlouisfed.org/series/FEDFUNDS
\item
  Bohn, H. (1998). The behavior of U.S. public debt and deficits.
  \emph{Quarterly Journal of Economics}, \emph{113}(3), 949-963.
  https://doi.org/10.1162/003355398555793
\item
  Bohn, H. (2005). \emph{The sustainability of fiscal policy in the
  United States} (CESifo Working Paper No.~1446). CESifo.
\item
  Bohn, H. (2007). Are stationarity and cointegration restrictions
  really necessary for the intertemporal budget constraint?
  \emph{Journal of Monetary Economics}, \emph{54}(7), 1837-1847.
  https://doi.org/10.1016/j.jmoneco.2006.12.012
\item
  Bohoslavsky, J. P., y Cantamutto, F. J. (2024). Deuda pública y
  desarrollo (no)sustentable en Argentina. Una visión crítica sobre la
  especialización productiva. \emph{FACES}, \emph{30}(62), 1-18.
  https://nulan.mdp.edu.ar/id/eprint/4109/1/FACES-30-62-bohoslavsky-cantamutto.pdf
\item
  Bouabdallah, O., Checherita-Westphal, C., Warmedinger, T., de Stefani,
  R., Drudi, F., Setzer, R., y Westphal, A. (2017). \emph{Debt
  sustainability analysis for euro area sovereigns: A methodological
  framework} (ECB Occasional Paper No.~185). Banco Central Europeo.
  https://ecb.europa.eu/pub/pdf/scpops/ecbop185.bg.pdf
\item
  Braun, O., y Joy, L. (1968). A model of economic stagnation: A case
  study of the Argentine economy. \emph{The Economic Journal},
  \emph{78}(312), 868-887. https://doi.org/10.2307/2229183
\item
  Calvo, G. A., Izquierdo, A., y Talvi, E. (2003). \emph{Sudden stops,
  the real exchange rate, and fiscal sustainability: Argentina's
  lessons} (NBER Working Paper No.~9828). National Bureau of Economic
  Research. https://doi.org/10.3386/w9828
\item
  Canova, F. (2005). The transmission of US shocks to Latin America.
  \emph{Journal of Applied Econometrics}, \emph{20}(2), 229-251.
  https://doi.org/10.1002/jae.837
\item
  Cantamutto, F. J., y Féliz, M. (2021). Argentina entre la
  sostenibilidad de la deuda y la vida. \emph{Revista Nuestramérica},
  \emph{9}(17), e6167845.
\item
  Cantamutto, F. J., Lobato, J., Rodríguez-Enríquez, C., y Féliz, M.
  (2024). Sustentabilidad de la deuda y derechos humanos desde la
  perspectiva de la sostenibilidad de la vida. En J. P. Bohoslavsky y M.
  L. Clérico (Eds.), \emph{Deuda y derechos humanos: Claves desde el
  sistema interamericano} (pp.~170-195). EDULP.
\item
  Celasun, O., Debrun, X., y Ostry, J. D. (2006). \emph{Primary surplus
  behavior and risks to fiscal sustainability in emerging market
  countries: A ``fan-chart'' approach} (IMF Working Paper No.~06/67).
  Fondo Monetario Internacional.
  https://doi.org/10.5089/9781451863277.001
\item
  Cetrángolo, O., Damill, M., Frenkel, R., y Jiménez, J. P. (1997).
  \emph{La sostenibilidad de la política fiscal en América Latina: El
  caso argentino} (Documento de Trabajo R-315). Banco Interamericano de
  Desarrollo, Red de Centros de Investigación.
  https://doi.org/10.18235/0011989
\item
  Christiano, L. J., Eichenbaum, M., y Evans, C. L. (1999). Monetary
  policy shocks: What have we learned and to what end? En J. B. Taylor y
  M. Woodford (Eds.), \emph{Handbook of Macroeconomics} (Vol. 1A,
  pp.~65-148). Elsevier.
\item
  Comín Comín, F. (2014). \emph{Historia económica mundial: De los
  orígenes a la actualidad}. Alianza Editorial.
\item
  Cushman, D. O., y Zha, T. (1997). Identifying monetary policy in a
  small open economy under flexible exchange rates. \emph{Journal of
  Monetary Economics}, \emph{39}(3), 433-448.
  https://doi.org/10.1016/S0304-3932(97)00029-9
\item
  Damill, M., Frenkel, R., y Rapetti, M. (2015). Macroeconomic policy in
  Argentina during 2002-2013. \emph{Comparative Economic Studies},
  \emph{57}(3), 369-400. https://doi.org/10.1057/ces.2015.3
\item
  della Paolera, G., y Taylor, A. M. (2001). \emph{Straining at the
  anchor: The Argentine currency board and the search for macroeconomic
  stability, 1880-1935}. University of Chicago Press.
\item
  Diamand, M. (1972). La estructura productiva desequilibrada argentina
  y el tipo de cambio. \emph{Desarrollo Económico}, \emph{12}(45),
  25-47. https://doi.org/10.2307/3465991
\item
  Díaz Alejandro, C. F. (1970). \emph{Essays on the economic history of
  the Argentine Republic}. Yale University Press.
\item
  Dickey, D. A., y Fuller, W. A. (1979). Distribution of the estimators
  for autoregressive time series with a unit root. \emph{Journal of the
  American Statistical Association}, \emph{74}(366), 427-431.
  https://doi.org/10.2307/2286348
\item
  Eichengreen, B., Hausmann, R., y Panizza, U. (2003). \emph{Currency
  mismatches, debt intolerance and original sin: Why they are not the
  same and why it matters} (NBER Working Paper No.~10036). National
  Bureau of Economic Research. https://doi.org/10.3386/w10036
\item
  Engle, R. F. (1982). Autoregressive conditional heteroscedasticity
  with estimates of the variance of United Kingdom inflation.
  \emph{Econometrica}, \emph{50}(4), 987-1007.
  https://doi.org/10.2307/1912773
\item
  Everaert, G., y Jansen, S. (2018). On the estimation of panel fiscal
  reaction functions: Heterogeneity or fiscal fatigue? \emph{Economic
  Modelling}, \emph{70}, 87-96.
  https://doi.org/10.1016/j.econmod.2017.10.014
\item
  Favero, C., y Giavazzi, F. (2007). \emph{Debt and the effects of
  fiscal policy} (NBER Working Paper No.~12822). National Bureau of
  Economic Research. https://doi.org/10.3386/w12822
\item
  Federal Reserve Bank of Atlanta. (s. f.). \emph{Wu-Xia shadow federal
  funds rate} {[}Conjunto de datos{]}.
  https://www.atlantafed.org/cqer/research/wu-xia-shadow-federal-funds-rate
\item
  Fondo Monetario Internacional {[}FMI{]}. (2021). \emph{Review of the
  debt sustainability framework for market access countries} (IMF Policy
  Paper).
\item
  Fondo Monetario Internacional {[}FMI{]}. (2024). \emph{Argentina:
  Eighth review under the extended arrangement under the Extended Fund
  Facility} (IMF Country Report No.~24/167).
\item
  Fondo Monetario Internacional {[}FMI{]}. (2025). \emph{Argentina:
  Request for an extended arrangement under the Extended Fund Facility}
  (IMF Country Report No.~25/95).
\item
  Fondo Monetario Internacional {[}FMI{]}. (2026). \emph{Argentina: 2026
  Article IV consultation, second review under the extended arrangement
  under the Extended Fund Facility} (IMF Country Report No.~26/105).
\item
  Fondo Monetario Internacional {[}FMI{]} y Mecanismo Europeo de
  Estabilidad {[}ESM{]}. (2023). \emph{The stochastic debt fanchart in
  the IMF's Sovereign Risk and Debt Sustainability Framework for MACs}
  {[}Presentación, 30 de noviembre de 2023{]}.
\item
  Ghosh, A. R., Kim, J. I., Mendoza, E. G., Ostry, J. D., y Qureshi, M.
  S. (2013). Fiscal fatigue, fiscal space and debt sustainability in
  advanced economies. \emph{The Economic Journal}, \emph{123}(566),
  F4-F30. https://doi.org/10.1111/ecoj.12010
\item
  Girouard, N., y André, C. (2005). \emph{Measuring cyclically-adjusted
  budget balances for OECD countries} (OECD Economics Department Working
  Papers No.~434). OECD Publishing.
\item
  Gonçalves, S., y Kilian, L. (2004). Bootstrapping autoregressions with
  conditional heteroskedasticity of unknown form. \emph{Journal of
  Econometrics}, \emph{123}(1), 89-120.
  https://doi.org/10.1016/j.jeconom.2003.10.030
\item
  González, M. J., y Hernández, J. M. (2023). \emph{Public debt
  sustainability and fiscal reaction functions in Latin America and the
  Caribbean} (IDB Discussion Paper IDB-DP-988). Banco Interamericano de
  Desarrollo. https://doi.org/10.18235/0004701
\item
  Hakkio, C. S., y Rush, M. (1991). Is the budget deficit ``too large?''
  \emph{Economic Inquiry}, \emph{29}(3), 429-445.
  https://doi.org/10.1111/j.1465-7295.1991.tb00837.x
\item
  Hébert, B., y Schreger, J. (2017). The costs of sovereign default:
  Evidence from Argentina. \emph{American Economic Review},
  \emph{107}(10), 3119-3145. https://doi.org/10.1257/aer.20151667
\item
  Hodrick, R. J., y Prescott, E. C. (1997). Postwar U.S. business
  cycles: An empirical investigation. \emph{Journal of Money, Credit and
  Banking}, \emph{29}(1), 1-16. https://doi.org/10.2307/2953682
\item
  Hosking, J. R. M. (1980). The multivariate portmanteau statistic.
  \emph{Journal of the American Statistical Association},
  \emph{75}(371), 602-608.
  https://doi.org/10.1080/01621459.1980.10477520
\item
  Infobae. (2026, 23 de septiembre). \emph{El riesgo país arriba de los
  550 puntos condiciona la primera emisión de deuda internacional de una
  provincia}.
  https://www.infobae.com/economia/2026/09/23/el-riesgo-pais-arriba-de-los-550-puntos-condiciona-la-primera-emision-de-deuda-internacional-de-una-provincia/
\item
  Ivanov, V., y Kilian, L. (2005). A practitioner's guide to lag order
  selection for VAR impulse response analysis. \emph{Studies in
  Nonlinear Dynamics and Econometrics}, \emph{9}(1), artículo 2.
  https://doi.org/10.2202/1558-3708.1219
\item
  Jarque, C. M., y Bera, A. K. (1987). A test for normality of
  observations and regression residuals. \emph{International Statistical
  Review}, \emph{55}(2), 163-172. https://doi.org/10.2307/1403192
\item
  Johansen, S. (1991). Estimation and hypothesis testing of
  cointegration vectors in Gaussian vector autoregressive models.
  \emph{Econometrica}, \emph{59}(6), 1551-1580.
  https://doi.org/10.2307/2938278
\item
  Johansen, S., Mosconi, R., y Nielsen, B. (2000). Cointegration
  analysis in the presence of structural breaks in the deterministic
  trend. \emph{The Econometrics Journal}, \emph{3}(2), 216-249.
  https://doi.org/10.1111/1368-423X.00047
\item
  Kalemli-Özcan, Ş. (2019). \emph{U.S. monetary policy and international
  risk spillovers} (NBER Working Paper No.~26297). National Bureau of
  Economic Research. https://doi.org/10.3386/w26297
\item
  Kaminsky, G., Lizondo, S., y Reinhart, C. M. (1998). Leading
  indicators of currency crises. \emph{IMF Staff Papers}, \emph{45}(1),
  1-48. https://doi.org/10.2307/3867328
\item
  Kilian, L. (1998). Small-sample confidence intervals for impulse
  response functions. \emph{Review of Economics and Statistics},
  \emph{80}(2), 218-230. https://doi.org/10.1162/003465398557465
\item
  Kwiatkowski, D., Phillips, P. C. B., Schmidt, P., y Shin, Y. (1992).
  Testing the null hypothesis of stationarity against the alternative of
  a unit root. \emph{Journal of Econometrics}, \emph{54}(1-3), 159-178.
  https://doi.org/10.1016/0304-4076(92)90104-Y
\item
  Leeper, E. M. (1991). Equilibria under ``active'' and ``passive''
  monetary and fiscal policies. \emph{Journal of Monetary Economics},
  \emph{27}(1), 129-147. https://doi.org/10.1016/0304-3932(91)90007-B
\item
  Levy Yeyati, E., y Sturzenegger, F. (2023). A balance-sheet approach
  to fiscal sustainability. \emph{Fiscal Studies}, \emph{44}(1), 61-84.
  https://doi.org/10.1111/1475-5890.12319
\item
  Ley 27.541 de Solidaridad Social y Reactivación Productiva en el Marco
  de la Emergencia Pública. (2019, 23 de diciembre). Boletín Oficial de
  la República Argentina.
\item
  Ley 27.612 de Fortalecimiento de la Sostenibilidad de la Deuda
  Pública. (2021, 3 de marzo). Boletín Oficial de la República
  Argentina.
  https://www.argentina.gob.ar/normativa/nacional/norma-347532/texto
\item
  Lütkepohl, H. (2005). \emph{New introduction to multiple time series
  analysis}. Springer.
\item
  Medeiros, J. (2012). \emph{Stochastic debt simulation using VAR models
  and a panel fiscal reaction function: Results for a selected number of
  countries} (European Economy Economic Papers No.~459). Comisión
  Europea. https://doi.org/10.2765/26084
\item
  Mendoza, E. G., y Oviedo, P. M. (2009). Public debt, fiscal solvency
  and macroeconomic uncertainty in Latin America: The cases of Brazil,
  Colombia, Costa Rica and Mexico. \emph{Economía Mexicana. Nueva
  Época}, \emph{18}(2), 133-177.
\item
  Ministerio de Economía de la Nación, Secretaría de Finanzas. (2025).
  \emph{Informe de deuda pública al 31/12/2025} {[}Conjunto de datos{]}.
\item
  Neumeyer, P. A., y Perri, F. (2005). Business cycles in emerging
  economies: The role of interest rates. \emph{Journal of Monetary
  Economics}, \emph{52}(2), 345-380.
  https://doi.org/10.1016/j.jmoneco.2004.04.011
\item
  Oficina de Presupuesto del Congreso {[}OPC{]}. (2026).
  \emph{Operaciones de deuda pública, julio 2026}.
\item
  Perron, P. (1989). The great crash, the oil price shock, and the unit
  root hypothesis. \emph{Econometrica}, \emph{57}(6), 1361-1401.
  https://doi.org/10.2307/1913712
\item
  Pesaran, M. H., y Smith, R. (1995). Estimating long-run relationships
  from dynamic heterogeneous panels. \emph{Journal of Econometrics},
  \emph{68}(1), 79-113. https://doi.org/10.1016/0304-4076(94)01644-F
\item
  Poder Ejecutivo Nacional. (2007). \emph{Decreto 897/2007: Fondo de
  Garantía de Sustentabilidad del Régimen Previsional Público de
  Reparto} (Boletín Oficial, 13 de julio de 2007). Base normativa de la
  AFIP. https://biblioteca.afip.gob.ar/dcp/DEC\_C\_000897\_2007\_07\_12
\item
  Porto, A. (2020). \emph{Evolución del sector público argentino en el
  largo plazo} (Documento de Trabajo N.° 36). CEFIP, Universidad
  Nacional de La Plata.
  https://www.mfp.econo.unlp.edu.ar/wp-content/uploads/2020/06/DocTrab36.pdf
\item
  Rapoport, M. (2020). \emph{Historia económica, política y social de la
  Argentina (1880-2003)}. Crítica.
\item
  Reinhart, C. M., y Rogoff, K. S. (2009). \emph{This time is different:
  Eight centuries of financial folly}. Princeton University Press.
\item
  Reinhart, C. M., Rogoff, K. S., y Savastano, M. A. (2003). Debt
  intolerance. \emph{Brookings Papers on Economic Activity},
  \emph{2003}(1), 1-74. https://doi.org/10.1353/eca.2003.0018
\item
  Renzhi, N., y Beirne, J. (2024). \emph{Global shocks and monetary
  policy transmission in emerging markets} (ADB Economics Working Paper
  No.~726). Asian Development Bank.
\item
  Rey, H. (2013). \emph{Dilemma not trilemma: The global financial cycle
  and monetary policy independence} {[}Ponencia{]}. Economic Policy
  Symposium, Federal Reserve Bank of Kansas City, Jackson Hole, Wyoming,
  Estados Unidos.
\item
  Rigobon, R. (2003). Identification through heteroskedasticity.
  \emph{The Review of Economics and Statistics}, \emph{85}(4), 777-792.
  https://doi.org/10.1162/003465303772815727
\item
  Rigobon, R., y Sack, B. (2004). The impact of monetary policy on asset
  prices. \emph{Journal of Monetary Economics}, \emph{51}(8), 1553-1575.
  https://doi.org/10.1016/j.jmoneco.2004.02.004
\item
  Rodríguez, A. (2023). Financiamiento y sostenibilidad de la deuda
  pública: El caso de la economía argentina (2001-2022). \emph{Ciencias
  Económicas}, \emph{2}(20), e0032.
  https://doi.org/10.14409/rce.2023.20.e0032
\item
  Roldán, V. (2021). \emph{La renegociación por la reestructuración de
  la deuda soberana argentina con los acreedores privados poseedores de
  instrumentos bajo jurisdicción extranjera durante 2020} {[}Tesina de
  grado{]}. Universidad Nacional de Rosario.
\item
  Sanches, P. (2019). \emph{Análisis de la sostenibilidad de la deuda
  soberana argentina. 2019-2023} {[}Tesis de maestría{]}. Universidad
  Torcuato Di Tella.
\item
  Sargent, T. J., y Wallace, N. (1981). Some unpleasant monetarist
  arithmetic. \emph{Federal Reserve Bank of Minneapolis Quarterly
  Review}, \emph{5}(3), 1-17.
\item
  Secretaría de Finanzas, Ministerio de Economía. (2025, 10 de
  diciembre). \emph{Resultado de licitación del Bono del Tesoro Nacional
  en dólares estadounidenses 6,50\% con vencimiento 30 de noviembre de
  2029} {[}Comunicado{]}.
  https://www.argentina.gob.ar/noticias/resultado-de-licitacion-del-bono-del-tesoro-nacional-en-dolares-estadounidenses-650-con
\item
  Sims, C. A., Stock, J. H., y Watson, M. W. (1990). Inference in linear
  time series models with some unit roots. \emph{Econometrica},
  \emph{58}(1), 113-144. https://doi.org/10.2307/2938337
\item
  Thirlwall, A. P. (1979). The balance of payments constraint as an
  explanation of international growth rate differences. \emph{Banca
  Nazionale del Lavoro Quarterly Review}, \emph{32}(128), 45-53.
\item
  Trehan, B., y Walsh, C. E. (1991). Testing intertemporal budget
  constraints: Theory and applications to U.S. federal budget and
  current account deficits. \emph{Journal of Money, Credit and Banking},
  \emph{23}(2), 206-223. https://doi.org/10.2307/1992777
\item
  Uribe, M., y Yue, V. Z. (2006). Country spreads and emerging
  countries: Who drives whom? \emph{Journal of International Economics},
  \emph{69}(1), 6-36. https://doi.org/10.1016/j.jinteco.2005.04.003
\item
  Wijayanti, R., y Rachmanira, S. (2020). Early warning system for
  government debt crisis in developing countries. \emph{Journal of
  Central Banking Theory and Practice}, \emph{9}(número especial 1),
  103-124. https://doi.org/10.2478/jcbtp-2020-0025
\item
  Wu, J. C., y Xia, F. D. (2016). Measuring the macroeconomic impact of
  monetary policy at the zero lower bound. \emph{Journal of Money,
  Credit and Banking}, \emph{48}(2-3), 253-291.
  https://doi.org/10.1111/jmcb.12300
\item
  Zivot, E., y Andrews, D. W. K. (1992). Further evidence on the great
  crash, the oil-price shock, and the unit-root hypothesis.
  \emph{Journal of Business and Economic Statistics}, \emph{10}(3),
  251-270. https://doi.org/10.1080/07350015.1992.10509904
\end{itemize}

\begin{center}\rule{0.5\linewidth}{0.5pt}\end{center}
